% Slides — Capítulo 5: Estabilidade Atmosférica
\documentclass[aspectratio=169,11pt]{beamer}
\usepackage{../comum/temameteo}
\graphicspath{{../figuras/cap05/}}

\title[Cap. 5 --- Estabilidade]{Capítulo 5 --- Estabilidade Atmosférica}
\subtitle{Meteorologia Prática}
\author{}
\date{}

\begin{document}

\begin{frame}
  \titlepage
  \vfill
  \atribuicao
\end{frame}

\begin{frame}{Roteiro}
  \tableofcontents
\end{frame}

% ---------------------------------------------------------------
\section{Diagramas termodinâmicos}

\begin{frame}{Construindo o diagrama: as famílias de linhas}
  \begin{columns}
    \column{0.25\textwidth}
    \begin{itemize}
      \item Isotermas/isóbaras;
      \item $+$ isoumas ($T_d$, LCL);
      \item $+$ adiabáticas secas ($\theta$);
      \item $+$ úmidas ($\theta_w$).
    \end{itemize}
    \column{0.37\textwidth}
    \figlivroslide{fig_05_1c.jpg}{0.4\textheight}{Fig.~5.1 ---
      adiabáticas secas no emagrama}
    \column{0.37\textwidth}
    \figlivroslide{fig_05_1d.jpg}{0.4\textheight}{Fig.~5.1 ---
      adiabáticas úmidas: mais suaves ($L_v$ devolvido)}
  \end{columns}
\end{frame}

\begin{frame}{O mesmo conteúdo, eixos deformados}
  \begin{columns}
    \column{0.4\textwidth}
    \begin{itemize}
      \item Emagrama, Skew-$T$, Stüve, tefigrama: mesmas 5 famílias.
      \item O Skew-$T$ \alert{inclina as isotermas} para abrir o ângulo
            com as adiabáticas: a sondagem fica legível.
      \item Padrão operacional mundial; o MetPy o desenha (Caso~1).
    \end{itemize}
    \column{0.3\textwidth}
    \figlivroslide{fig_05_3b.jpg}{0.42\textheight}{Fig.~5.3 --- emagrama}
    \column{0.3\textwidth}
    \figlivroslide{fig_05_3c.jpg}{0.42\textheight}{Fig.~5.3 ---
      Skew-$T$ log-$P$}
  \end{columns}
\end{frame}

\begin{frame}{O Skew-$T$ completo}
  \begin{columns}
    \column{0.44\textwidth}
    \begin{itemize}
      \item Sondagem: $T$ e $T_d$ ponto a ponto \alert{sobre} as réguas.
      \item Parcela de superfície: seca até o LCL (Normand, Cap.~4),
            úmida acima.
      \item Áreas entre parcela e ambiente $=$ energia (CAPE/CIN,
            Cap.~14).
    \end{itemize}
    \column{0.56\textwidth}
    \figlivroslide{fig_05_6g.jpg}{0.58\textheight}{Fig.~5.6 --- as
      famílias sobrepostas no Skew-$T$ log-$P$}
  \end{columns}
\end{frame}

\begin{frame}{Processos no diagrama: subida até a nuvem}
  \begin{columns}
    \column{0.4\textwidth}
    \begin{itemize}
      \item Subida seca: $T$ pela adiabática, $T_d$ pela isouma — as
            duas se encontram no \alert{LCL}.
      \item Acima: adiabática úmida (condensa e aquece).
      \item Descida: seca de novo (a chuva ficou lá em cima) — föhn
            (Cap.~3, Cap.~17).
    \end{itemize}
    \column{0.6\textwidth}
    \figlivroslide{fig_05_9.jpg}{0.56\textheight}{Fig.~5.9 --- parcela
      levantada: LCL como cruzamento adiabática seca × isouma}
  \end{columns}
\end{frame}

% ---------------------------------------------------------------
\section{Método da parcela e empuxo}

\begin{frame}{Quem sobe e quem afunda: $T_v$ decide}
  \begin{columns}
    \column{0.5\textwidth}
    \eqdestaque{$\dfrac{F}{m} = |g|\,
      \dfrac{T_{v,p}-T_{v,e}}{T_{v,e}}$}
    \vspace{0.4em}
    \begin{itemize}
      \item Arquimedes $+$ lei dos gases; \alert{virtual} porque vapor
            alivia e gotículas pesam (Cap.~1).
      \item $2$\,K de excesso $\Rightarrow$ $0{,}065$\,m/s$^2$
            $\Rightarrow$ $11$\,m/s após 1\,km: um termal de planador.
      \item 1\,g/kg de água líquida $\simeq -0{,}3$\,K (T1).
    \end{itemize}
    \column{0.5\textwidth}
    \figlivroslide{fig_05_12c.jpg}{0.5\textheight}{Fig.~5.12 --- parcela
      × ambiente: o processo adiabático contra o perfil}
  \end{columns}
\end{frame}

\begin{frame}{A atmosfera estável é uma mola}
  \begin{columns}
    \column{0.44\textwidth}
    \eqdestaque{$\ddot z' = -N^2 z'$,\qquad
      $N^2 = \dfrac{|g|}{\theta_v}\dfrac{\Delta\theta_v}{\Delta z}$}
    \vspace{0.4em}
    \begin{itemize}
      \item Troposfera: $N \simeq 0{,}01$\,s$^{-1}$,
            $P_{BV} \simeq 10$\,min.
      \item Visível: ondas de montanha, faixas de nuvens espaçadas de
            $U P_{BV}$ (Cap.~17).
      \item Simulação com \texttt{solve\_ivp}: Caso~2.
    \end{itemize}
    \column{0.56\textwidth}
    \figlivroslide{fig_05_13c.jpg}{0.5\textheight}{Fig.~5.13 --- a
      analogia da mola: oscilação de Brunt--Väisälä}
  \end{columns}
\end{frame}

% ---------------------------------------------------------------
\section{Estabilidade estática: local e não local}

\begin{frame}{As cinco classes de estabilidade}
  \begin{columns}
    \column{0.3\textwidth}
    \begin{itemize}
      \item Comparar $\gamma$ do ambiente com $\Gamma_d$ e $\Gamma_s$.
      \item \alert{Condicional}: estável seca, instável saturada — a
            atmosfera ``armada'' dos trópicos (Cap.~14).
    \end{itemize}
    \column{0.35\textwidth}
    \figlivroslide{fig_05_15a.jpg}{0.44\textheight}{Fig.~5.15 ---
      absolutamente estável e saturada neutra}
    \column{0.35\textwidth}
    \figlivroslide{fig_05_15b.jpg}{0.44\textheight}{Fig.~5.15 ---
      condicionalmente instável}
  \end{columns}
\end{frame}

\begin{frame}{O critério local falha: estabilidade não local}
  \begin{columns}
    \column{0.4\textwidth}
    \begin{itemize}
      \item Local: sinal de $\Delta\theta_v/\Delta z$ no nível.
      \item \alert{Não local}: siga as \emph{trajetórias} — camada
            ``neutra'' atravessada por parcela quente de baixo é
            instável.
      \item Método do Stull: levantar/afundar parcelas de todos os
            níveis (Caso~3 do notebook).
    \end{itemize}
    \column{0.6\textwidth}
    \figlivroslide{fig_05_14cb.jpg}{0.56\textheight}{Fig.~5.14 ---
      classificação não local sobre a sondagem: zonas turbulentas}
  \end{columns}
\end{frame}

% ---------------------------------------------------------------
\section{Estabilidade dinâmica: Richardson}

\begin{frame}{Cisalhamento contra a mola}
  \begin{columns}
    \column{0.44\textwidth}
    \eqdestaque{$Ri = \dfrac{N^2}{|\Delta\vec U/\Delta z|^2}$;\qquad
      $Ri < 0{,}25 \Rightarrow$ turbulência}
    \vspace{0.4em}
    \begin{itemize}
      \item Ar estável PODE virar turbulento: basta cisalhamento.
      \item Kelvin--Helmholtz: a instabilidade fotografável.
      \item CAT no jato (Cap.~11); camada noturna intermitente
            (Cap.~18). $Ri(z)$ da sondagem: Caso~4.
    \end{itemize}
    \column{0.56\textwidth}
    \figlivroslide{fig_05_17.jpg}{0.52\textheight}{Fig.~5.17 --- ondas de
      Kelvin--Helmholtz nascendo do cisalhamento}
  \end{columns}
\end{frame}

\begin{frame}{Richardson na prática}
  \begin{columns}
    \column{0.44\textwidth}
    \begin{itemize}
      \item Perfis medidos de $U$, $V$, $\theta$ $\to$ $Ri(z)$ camada a
            camada.
      \item $Ri$ atravessa 0{,}25 nas zonas de cisalhamento do jato
            noturno.
      \item Exercício T4: o jato com $\Delta U = 20$\,m/s em 500\,m.
    \end{itemize}
    \column{0.56\textwidth}
    \figlivroslide{fig_05_17a.jpg}{0.52\textheight}{Fig.~5.17 --- perfis
      de vento e $\theta$ e o $Ri$ resultante}
  \end{columns}
\end{frame}

% ---------------------------------------------------------------
\section{Tropopausa e estrutura do perfil}

\begin{frame}{Lendo a sondagem inteira}
  \begin{columns}
    \column{0.4\textwidth}
    \begin{itemize}
      \item \alert{Camada de mistura}: $\theta_v$ constante até $z_i$.
      \item \alert{Inversões}: tampas de poluição (Cap.~19) e de
            convecção (Cap.~14).
      \item \alert{Tropopausa}: salto de $N^2$; WMO: gradiente
            $\le2$\,K/km. A tampa global.
    \end{itemize}
    \column{0.6\textwidth}
    \figlivroslide{fig_05_18.jpg}{0.56\textheight}{Fig.~5.18 --- camada
      de mistura, camadas isotérmicas, inversão e tropopausa no
      Skew-$T$}
  \end{columns}
\end{frame}

\begin{frame}{A tropopausa no planeta}
  \begin{columns}
    \column{0.44\textwidth}
    \begin{itemize}
      \item Alta e fria nos trópicos ($\sim16$\,km); baixa nos polos
            ($\sim8$\,km); degrau perto do jato (Cap.~11).
      \item Isentrópicas ($\theta$ const) cruzam a tropopausa nas
            dobras — intercâmbio estratosfera--troposfera.
      \item Bigornas de tempestade se achatam nela (Cap.~14).
    \end{itemize}
    \column{0.56\textwidth}
    \figlivroslide{fig_05_20.jpg}{0.52\textheight}{Fig.~5.20 ---
      tropopausa e isentrópicas em corte meridional}
  \end{columns}
\end{frame}

% ---------------------------------------------------------------
\section{Síntese e laboratório}

\begin{frame}{O mapa do capítulo}
  \eqdestaque{$\dfrac{F}{m} = |g|\dfrac{T_{v,p}-T_{v,e}}{T_{v,e}}
    \;\to\; N^2 = \dfrac{|g|}{\theta_v}\dfrac{\Delta\theta_v}{\Delta z}
    \;\to\; \text{não local!} \;\to\; Ri \lessgtr 0{,}25$}
  \vspace{0.6em}
  Empuxo: volta ou dispara. $N$: com que ritmo volta. $Ri$: o
  cisalhamento bagunça mesmo quando voltaria. Tabuleiro: o Skew-$T$.
\end{frame}

\begin{frame}{Laboratório numérico --- notebook do Cap.~5}
  \texttt{notebooks/cap05\_estabilidade.ipynb}
  \begin{enumerate}
    \item Skew-$T$ do zero e com MetPy (sondagem de Campo de Marte).
    \item Oscilador de Brunt--Väisälä com \texttt{solve\_ivp}.
    \item Estabilidade não local: o algoritmo do Stull em NumPy.
    \item $Ri(z)$ e o mapa de turbulência da sondagem.
  \end{enumerate}
  \vspace{0.4em}
  \textbf{Exercícios}: T1--T5 e N1--N4 nas notas; células reservadas no
  notebook.
  \vfill
  \atribuicao
\end{frame}

\end{document}
